\section{Formal glueing of module categories}
\label{section-formal-glueing}

\noindent
Fix a Noetherian scheme $X$, and a closed subscheme $Z$ with complement $U$.
Our goal is to explain how coherent sheaves on $X$ can be constructed
(uniquely) from coherent sheaves on the formal completion of $X$ along
$Z$, and those on $U$ with a suitable compatibility on the overlap.
We first do this using only commutative algebra (this section) and
later we explain this in the setting of algebraic spaces
(Pushouts of Spaces, Section \ref{spaces-pushouts-section-formal-glueing}).

\medskip\noindent
Here are some references treating some of the material in this section:
\cite[Section 2]{ArtinII},
\cite[Appendix]{Ferrand-Raynaud},
\cite{Beauville-Laszlo},
\cite{MB}, and
\cite[Section 4.6]{dJ-crystalline}.

\begin{lemma}
\label{lemma-characterize-flatness-on-torsion}
Let $\varphi : R \to S$ be a ring map. Let $I \subset R$ be an ideal.
The following are equivalent
\begin{enumerate}
\item $\varphi$ is flat and $R/I \to S/IS$ is faithfully flat,
\item $\varphi$ is flat, and the map
$\Spec(S/IS) \to \Spec(R/I)$ is surjective.
\item $\varphi$ is flat, and the base change functor
$M \mapsto M \otimes_R S$ is faithful on modules annihilated by $I$, and
\item $\varphi$ is flat, and the base change functor
$M \mapsto M \otimes_R S$ is faithful on $I$-power torsion modules.
\end{enumerate}
\end{lemma}

\begin{proof}
If $R \to S$ is flat, then $R/I^n \to S/I^nS$ is flat for every $n$, see
Algebra, Lemma \ref{algebra-lemma-flat-base-change}.
Hence (1) and (2) are equivalent by
Algebra, Lemma \ref{algebra-lemma-ff-rings}.
The equivalence of (1) with (3) follows by identifying $I$-torsion
$R$-modules with $R/I$-modules, using that
$$
M \otimes_R S = M \otimes_{R/I} S/IS
$$
for $R$-modules $M$ annihilated by $I$, and
Algebra, Lemma \ref{algebra-lemma-easy-ff}.
The implication (4) $\Rightarrow$ (3) is immediate. Assume (3). We have
seen above that $R/I^n \to S/I^nS$ is flat, and by assumption it induces
a surjection on spectra, as $\Spec(R/I^n) = \Spec(R/I)$ and
similarly for $S$. Hence the base change functor is faithful on modules
annihilated by $I^n$. Since any $I$-power torsion module $M$ is the union
$M = \bigcup M_n$ where $M_n$ is annihilated by $I^n$ we see that the base
change functor is faithful on the category of all $I$-power torsion modules
(as tensor product commutes with colimits).
\end{proof}

\begin{lemma}
\label{lemma-neighbourhood-isomorphism}
Assume $(\varphi : R \to S, I)$ satisfies the equivalent conditions of
Lemma \ref{lemma-characterize-flatness-on-torsion}.
The following are equivalent
\begin{enumerate}
\item for any $I$-power torsion module $M$, the natural map
$M \to M \otimes_R S$ is an isomorphism, and
\item $R/I \to S/IS$ is an isomorphism.
\end{enumerate}
\end{lemma}

\begin{proof}
The implication (1) $\Rightarrow$ (2) is immediate.
Assume (2). First assume that $M$ is annihilated by $I$.
In this case, $M$ is an $R/I$-module. Hence, we have an isomorphism
$$
M \otimes_R S = M \otimes_{R/I} S/IS = M \otimes_{R/I} R/I = M
$$
proving the claim. Next we prove by induction that $M \to M \otimes_R S$
is an isomorphism for any module $M$ is annihilated by $I^n$. Assume
the induction hypothesis holds for $n$ and assume $M$ is annihilated by
$I^{n + 1}$. Then we have a short exact sequence
$$
0 \to I^nM \to M \to M/I^nM \to 0
$$
and as $R \to S$ is flat this gives rise to a short exact sequence
$$
0 \to I^nM \otimes_R S \to M \otimes_R S \to M/I^nM \otimes_R S \to 0
$$
Using that the canonical map is an isomorphism for $M' = I^nM$ and
$M'' = M/I^nM$ (by induction hypothesis) we conclude the same thing is
true for $M$. Finally, suppose that $M$ is a general $I$-power torsion
module. Then $M = \bigcup M_n$ where $M_n$ is annihilated by $I^n$
and we conclude using that tensor products commute with colimits.
\end{proof}

\begin{lemma}
\label{lemma-neighbourhood-equivalence}
Assume $\varphi : R \to S$ is a flat ring map and $I \subset R$ is a
finitely generated ideal such that $R/I \to S/IS$ is an isomorphism. Then
\begin{enumerate}
\item for any $R$-module $M$ the map $M \to M \otimes_R S$ induces
an isomorphism
$M[I^\infty] \to (M \otimes_R S)[(IS)^\infty]$ of $I$-power
torsion submodules,
\item the natural map
$$
\Hom_R(M, N) \longrightarrow \Hom_S(M \otimes_R S, N \otimes_R S)
$$
is an isomorphism if either $M$ or $N$ is $I$-power torsion, and
\item the base change functor $M \mapsto M \otimes_R S$ defines an
equivalence of categories between $I$-power torsion modules
and $IS$-power torsion modules.
\end{enumerate}
\end{lemma}

\begin{proof}
Note that the equivalent conditions of both
Lemma \ref{lemma-characterize-flatness-on-torsion} and
Lemma \ref{lemma-neighbourhood-isomorphism}
are satisfied. We will use these without further mention.
We first prove (1). Let $M$ be any $R$-module.
Set $M' = M/M[I^\infty]$ and consider the exact sequence
$$
0 \to M[I^\infty] \to M \to M' \to 0
$$
As $M[I^\infty] = M[I^\infty] \otimes_R S$ we see that it suffices to
show that $(M' \otimes_R S)[(IS)^\infty] = 0$.
Write $I = (f_1, \ldots, f_t)$. By
Lemma \ref{lemma-divide-by-torsion}
we see that $M'[I^\infty] = 0$. Hence for every $n > 0$ the map
$$
M' \longrightarrow \bigoplus\nolimits_{i = 1, \ldots t} M',
\quad
x \longmapsto (f_1^n x, \ldots, f_t^n x)
$$
is injective. As $S$ is flat over $R$ also the corresponding map
$M' \otimes_R S \to \bigoplus_{i = 1, \ldots t} M' \otimes_R S$
is injective. This means that $(M' \otimes_R S)[I^n] = 0$ as desired.

\medskip\noindent
Next we prove (2). If $N$ is $I$-power torsion, then
$N \otimes_R S = N$ and the displayed map of (2) is an isomorphism by
Algebra, Lemma \ref{algebra-lemma-adjoint-tensor-restrict}.
If $M$ is $I$-power torsion, then the image of any map
$M \to N$ factors through $N[I^\infty]$ and the image of any map
$M \otimes_R S \to N \otimes_R S$ factors through
$(N \otimes_R S)[(IS)^\infty]$. Hence in this case
part (1) guarantees that we may replace $N$ by $N[I^\infty]$
and the result follows from the case where $N$ is $I$-power torsion
we just discussed.

\medskip\noindent
Next we prove (3). The functor is fully faithful by (2).
For essential surjectivity, we simply note that for any $IS$-power torsion
$S$-module $N$, the natural map $N \otimes_R S \to N$ is an isomorphism.
\end{proof}

\begin{lemma}
\label{lemma-map-identifies-koszul-and-cech-complexes}
Assume $\varphi : R \to S$ is a flat ring map and $I \subset R$ is a
finitely generated ideal such that $R/I \to S/IS$ is an isomorphism.
For any $f_1, \ldots, f_r \in R$ such that $V(f_1, \ldots, f_r) = V(I)$
\begin{enumerate}
\item the map of Koszul complexes
$K(R, f_1, \ldots, f_r) \to K(S, f_1, \ldots, f_r)$ is a quasi-isomorphism, and
\item The map of extended alternating {\v C}ech complexes
$$
\xymatrix{
R \to \prod_{i_0} R_{f_{i_0}} \to \prod_{i_0 < i_1} R_{f_{i_0}f_{i_1}}
\to \ldots \to R_{f_1\ldots f_r} \ar[d] \\
S \to \prod_{i_0} S_{f_{i_0}} \to \prod_{i_0 < i_1} S_{f_{i_0}f_{i_1}}
\to \ldots \to S_{f_1\ldots f_r}
}
$$
is a quasi-isomorphism.
\end{enumerate}
\end{lemma}

\begin{proof}
In both cases we have a complex $K_\bullet$ of $R$ modules and we want
to show that $K_\bullet \to K_\bullet \otimes_R S$ is a quasi-isomorphism.
By Lemma \ref{lemma-neighbourhood-isomorphism} and the flatness of
$R \to S$ this will hold as soon as all homology groups of $K$ are $I$-power
torsion. This is true for the Koszul complex by
Lemma \ref{lemma-homotopy-koszul} and for the extended alternating {\v C}ech
complex by Lemma \ref{lemma-extended-alternating-torsion}.
\end{proof}

\begin{lemma}
\label{lemma-naive-Koszul-complex}
Let $R$ be a ring. Let $I = (f_1, \ldots, f_n)$ be a finitely generated ideal
of $R$. Let $M$ be the $R$-module generated by elements
$e_1, \ldots, e_n$ subject to the relations $f_i e_j - f_j e_i = 0$.
There exists a short exact sequence
$$
0 \to K \to M \to I \to 0
$$
such that $K$ is annihilated by $I$.
\end{lemma}

\begin{proof}
This is just a truncation of the Koszul complex.
The map $M \to I$ is determined by the rule $e_i \mapsto f_i$. If
$m = \sum a_i e_i$ is in the kernel of $M \to I$, i.e., $\sum a_i f_i = 0$,
then $f_j m = \sum f_j a_i e_i = (\sum f_i a_i) e_j = 0$.
\end{proof}

\begin{lemma}
\label{lemma-explicit-ext}
Let $R$ be a ring. Let $I = (f_1, \ldots, f_n)$ be a finitely generated ideal
of $R$. For any $R$-module $N$ set
$$
H_1(N, f_\bullet) =
\frac{\{(x_1, \ldots, x_n) \in N^{\oplus n} \mid f_i x_j = f_j x_i \}}
{\{f_1x, \ldots, f_nx) \mid x \in N\}}
$$
For any $R$-module $N$ there exists a canonical short exact sequence
$$
0 \to \Ext_R(R/I, N) \to H_1(N, f_\bullet) \to \Hom_R(K, N)
$$
where $K$ is as in
Lemma \ref{lemma-naive-Koszul-complex}.
\end{lemma}

\begin{proof}
The notation above indicates the $\Ext$-groups in $\text{Mod}_R$
as defined in
Homology, Section \ref{homology-section-extensions}.
These are denoted $\Ext_R(M, N)$. Using the long exact sequence of
Homology, Lemma \ref{homology-lemma-six-term-sequence-ext}
associated to the short exact sequence $0 \to I \to R \to R/I \to 0$
and the fact that $\Ext_R(R, N) = 0$ we see that
$$
\Ext_R(R/I, N) =
\Coker(N \longrightarrow \Hom(I, N))
$$
Using the short exact sequence of
Lemma \ref{lemma-naive-Koszul-complex}
we see that we get a complex
$$
N \to \Hom(M, N) \to \Hom_R(K, N)
$$
whose homology in the middle is canonically isomorphic to
$\Ext_R(R/I, N)$. The proof of the lemma is now complete
as the cokernel of the first map
is canonically isomorphic to $H_1(N, f_\bullet)$.
\end{proof}

\begin{lemma}
\label{lemma-koszul-homology-annihilated}
Let $R$ be a ring. Let $I = (f_1, \ldots, f_n)$ be a finitely generated ideal
of $R$. For any $R$-module $N$ the Koszul homology group
$H_1(N, f_\bullet)$ defined in
Lemma \ref{lemma-explicit-ext}
is annihilated by $I$.
\end{lemma}

\begin{proof}
Let $(x_1, \ldots, x_n) \in N^{\oplus n}$ with $f_i x_j = f_j x_i$.
Then we have $f_i(x_1, \ldots, x_n) = (f_i x_i, \ldots, f_i x_n)$.
In other words $f_i$ annihilates $H_1(N, f_\bullet)$.
\end{proof}

\noindent
We can improve on the full faithfulness of
Lemma \ref{lemma-neighbourhood-equivalence}
by showing that $\Ext$-groups whose source is $I$-power torsion
are insensitive to passing to $S$ as well. See
Dualizing Complexes, Lemma \ref{dualizing-lemma-neighbourhood-extensions}
for a derived version of the following lemma.

\begin{lemma}
\label{lemma-neighbourhood-extensions}
Assume $\varphi : R \to S$ is a flat ring map and $I \subset R$ is a
finitely generated ideal such that $R/I \to S/IS$ is an isomorphism.
Let $M$, $N$ be $R$-modules. Assume $M$ is $I$-power torsion.
Given an short exact sequence
$$
0 \to N \otimes_R S \to \tilde E \to M \otimes_R S \to 0
$$
there exists a commutative diagram
$$
\xymatrix{
0 \ar[r] &
N \ar[r] \ar[d] &
E \ar[r] \ar[d] &
M \ar[r] \ar[d] &
0 \\
0 \ar[r] &
N \otimes_R S \ar[r] &
\tilde E \ar[r] &
M \otimes_R S \ar[r] &
0
}
$$
with exact rows.
\end{lemma}

\begin{proof}
As $M$ is $I$-power torsion we see that $M \otimes_R S = M$, see
Lemma \ref{lemma-neighbourhood-isomorphism}.
We will use this identification without further mention.
As $R \to S$ is flat, the base change functor is exact and we
obtain a functorial map of $\Ext$-groups
$$
\Ext_R(M, N)
\longrightarrow
\Ext_S(M \otimes_R S, N \otimes_R S),
$$
see
Homology, Lemma \ref{homology-lemma-exact-functor-ext}.
The claim of the lemma is that this map is surjective when
$M$ is $I$-power torsion. In fact we will show that it is an
isomorphism. By
Lemma \ref{lemma-I-power-torsion-presentation}
we can find a surjection $M' \to M$ with $M'$ a direct sum of
modules of the form $R/I^n$. Using the long exact sequence of
Homology, Lemma \ref{homology-lemma-six-term-sequence-ext}
and using Lemma \ref{lemma-neighbourhood-equivalence}
we see that it suffices to prove the lemma for $M'$.
Using compatibility of $\Ext$ with direct sums (details omitted)
we reduce to the case where $M = R/I^n$ for some $n$.

\medskip\noindent
Let $f_1, \ldots, f_t$ be generators for $I^n$. By
Lemma \ref{lemma-explicit-ext}
we have a commutative diagram
$$
\xymatrix{
0 \ar[r] &
\Ext_R(R/I^n, N) \ar[r] \ar[d] &
H_1(N, f_\bullet) \ar[r] \ar[d] &
\Hom_R(K, N) \ar[d] \\
0 \ar[r] &
\Ext_S(S/I^nS, N \otimes S) \ar[r] &
H_1(N \otimes S, f_\bullet) \ar[r] &
\Hom_S(K \otimes S, N \otimes S)
}
$$
with exact rows where $K$ is as in
Lemma \ref{lemma-naive-Koszul-complex}.
Hence it suffices to prove that the two right vertical arrows are
isomorphisms. Since $K$ is annihilated by $I^n$ we see that
$\Hom_R(K, N) = \Hom_S(K \otimes_R S, N \otimes_R S)$ by
Lemma \ref{lemma-neighbourhood-equivalence}.
As $R \to S$ is flat we have
$H_1(N, f_\bullet) \otimes_R S = H_1(N \otimes_R S, f_\bullet)$.
As $H_1(N, f_\bullet)$ is annihilated by $I^n$, see
Lemma \ref{lemma-koszul-homology-annihilated}
we have $H_1(N, f_\bullet) \otimes_R S = H_1(N, f_\bullet)$ by
Lemma \ref{lemma-neighbourhood-isomorphism}.
\end{proof}

\noindent
Let $R \to S$ be a ring map.
Let $f_1, \ldots, f_t \in R$ and $I = (f_1, \ldots, f_t)$.
Then for any $R$-module $M$ we can define a complex
\begin{equation}
\label{equation-glueing-complex}
0 \to M \xrightarrow{\alpha}
M \otimes_R S \times \prod M_{f_i} \xrightarrow{\beta}
\prod (M \otimes_R S)_{f_i}
\times
\prod M_{f_if_j}
\end{equation}
where $\alpha(m) = (m \otimes 1, m/1, \ldots, m/1)$ and
$$
\beta(m', m_1, \ldots, m_t) =
((m'/1 - m_1 \otimes 1, \ldots, m'/1 - m_t \otimes 1),
(m_1 - m_2, \ldots, m_{t - 1} - m_t).
$$
We would like to know when this complex is exact.

\begin{lemma}
\label{lemma-recover-module-from-glueing-data}
Assume $\varphi : R \to S$ is a flat ring map and
$I = (f_1, \ldots, f_t) \subset R$ is an ideal such that
$R/I \to S/IS$ is an isomorphism.
Let $M$ be an $R$-module. Then the
complex (\ref{equation-glueing-complex})
is exact.
\end{lemma}

\begin{proof}
First proof. Denote $\check{\mathcal{C}}_R \to \check{\mathcal{C}}_S$
the quasi-isomorphism of extended alternating {\v C}ech complexes of
Lemma \ref{lemma-map-identifies-koszul-and-cech-complexes}.
Since these complexes are bounded with flat terms, we see that
$M \otimes_R \check{\mathcal{C}}_R \to M \otimes_R \check{\mathcal{C}}_S$
is a quasi-isomorphism too (Lemmas
\ref{lemma-derived-tor-quasi-isomorphism} and
\ref{lemma-derived-tor-quasi-isomorphism-other-side}). Now the complex
(\ref{equation-glueing-complex}) is a truncation of the cone
of the map
$M \otimes_R \check{\mathcal{C}}_R \to M \otimes_R \check{\mathcal{C}}_S$
and we win.

\medskip\noindent
Second computational proof.
Let $m \in M$. If $\alpha(m) = 0$, then $m \in M[I^\infty]$, see
Lemma \ref{lemma-torsion-free}. Pick $n$ such that $I^n m = 0$
and consider the map $\varphi : R/I^n \to M$.
If $m \otimes 1 = 0$, then $\varphi \otimes 1_S = 0$, hence
$\varphi = 0$ (see
Lemma \ref{lemma-neighbourhood-equivalence})
hence $m = 0$. In this way we see that $\alpha$ is injective.

\medskip\noindent
Let $(m', m'_1, \ldots, m'_t) \in \Ker(\beta)$.
Write $m'_i = m_i/f_i^n$ for some $n > 0$ and $m_i \in M$.
We may, after possibly enlarging $n$ assume that
$f_i^n m' = m_i \otimes 1$ in $M \otimes_R S$ and
$f_j^nm_i - f_i^nm_j = 0$ in $M$.
In particular we see that
$(m_1, \ldots, m_t)$ defines an element $\xi$ of
$H_1(M, (f_1^n, \ldots, f_t^n))$.
Since $H_1(M, (f_1^n, \ldots, f_t^n))$ is annihilated by $I^{tn + 1}$ (see
Lemma \ref{lemma-koszul-homology-annihilated})
and since $R \to S$ is flat we see that
$$
H_1(M, (f_1^n, \ldots, f_t^n)) =
H_1(M, (f_1^n, \ldots, f_t^n)) \otimes_R S =
H_1(M \otimes_R S, (f_1^n, \ldots, f_t^n))
$$
by
Lemma \ref{lemma-neighbourhood-isomorphism}
The existence of $m'$ implies that $\xi$ maps to zero in the last group, i.e.,
the element $\xi$ is zero. Thus there exists an $m \in M$ such that
$m_i = f_i^n m$. Then $(m', m'_1, \ldots, m'_t) - \alpha(m)
= (m'', 0, \ldots, 0)$ for some $m'' \in (M \otimes_R S)[(IS)^\infty]$.
By
Lemma \ref{lemma-neighbourhood-equivalence}
we conclude that $m'' \in M[I^\infty]$ and we win.
\end{proof}

\begin{remark}
\label{remark-glueing-data}
In this remark we define a category of glueing data.
Let $R \to S$ be a ring map.
Let $f_1, \ldots, f_t \in R$ and $I = (f_1, \ldots, f_t)$.
Consider the category $\text{Glue}(R \to S, f_1, \ldots, f_t)$
as the category whose
\begin{enumerate}
\item objects are systems $(M', M_i, \alpha_i, \alpha_{ij})$, where
$M'$ is an $S$-module, $M_i$ is an $R_{f_i}$-module,
$\alpha_i : (M')_{f_i} \to M_i \otimes_R S$ is an isomorphism, and
$\alpha_{ij} : (M_i)_{f_j} \to (M_j)_{f_i}$ are isomorphisms
such that
\begin{enumerate}
\item $\alpha_{ij} \circ \alpha_i = \alpha_j$ as maps
$(M')_{f_if_j} \to (M_j)_{f_i}$, and
\item $\alpha_{jk} \circ \alpha_{ij} = \alpha_{ik}$ as maps
$(M_i)_{f_jf_k} \to (M_k)_{f_if_j}$ (cocycle condition).
\end{enumerate}
\item morphisms
$(M', M_i, \alpha_i, \alpha_{ij}) \to (N', N_i, \beta_i, \beta_{ij})$
are given by maps $\varphi' : M' \to N'$ and $\varphi_i : M_i \to N_i$
compatible with the given maps $\alpha_i, \beta_i, \alpha_{ij}, \beta_{ij}$.
\end{enumerate}
There is a canonical functor
$$
\text{Can} : \text{Mod}_R
\longrightarrow
\text{Glue}(R \to S, f_1, \ldots, f_t),
\quad
M \longmapsto (M \otimes_R S, M_{f_i}, \text{can}_i, \text{can}_{ij})
$$
where $\text{can}_i : (M \otimes_R S)_{f_i} \to M_{f_i} \otimes_R S$
and $\text{can}_{ij} : (M_{f_i})_{f_j} \to (M_{f_j})_{f_i}$
are the canonical isomorphisms. For any object
$\mathbf{M} = (M', M_i, \alpha_i, \alpha_{ij})$ of the category
$\text{Glue}(R \to S, f_1, \ldots, f_t)$ we define
$$
H^0(\mathbf{M}) =
\{(m', m_i) \mid \alpha_i(m') = m_i \otimes 1, \alpha_{ij}(m_i) = m_j\}
$$
in other words defined by the exact sequence
$$
0 \to H^0(\mathbf{M}) \to
M' \times \prod M_i \to
\prod M'_{f_i}
\times
\prod (M_i)_{f_j}
$$
similar to (\ref{equation-glueing-complex}).
We think of $H^0(\mathbf{M})$ as an $R$-module. Thus we also get a functor
$$
H^0 :
\text{Glue}(R \to S, f_1, \ldots, f_t)
\longrightarrow
\text{Mod}_R
$$
Our next goal is to show that the functors
$\text{Can}$ and $H^0$ are sometimes quasi-inverse to each other.
\end{remark}

\begin{lemma}
\label{lemma-H0-Can-adjoint}
In Remark \ref{remark-glueing-data} the functor
$H^0 : \text{Glue}(R \to S, f_1, \ldots, f_t) \to \text{Mod}_R$
is a right adjoint to the functor
$\text{Can} : \text{Mod}_R \to \text{Glue}(R \to S, f_1, \ldots, f_t)$.
\end{lemma}

\begin{proof}
Let $\mathbf{M} = (M', M_i, \alpha_i, \alpha_{ij})$ be an object of
$\text{Glue}(R \to S, f_1, \ldots, f_t)$. For any $R$-module $N$ there is
a map
$$
\Hom_{\text{Glue}(R \to S, f_1, \ldots, f_t)}(\text{Can}(N), \mathbf{M})
\to
\Hom_R(N, H^0(\mathbf{M}))
$$
sending $\psi$ to $H^0(\psi)$ composed with the obvious map
$N \to H^0(\text{Can}(N))$. By construction the displayed map is an isomorphism
for $N = R$ (even if $R \to H^0(\text{Can}(R))$ is not an isomorphism
in general). The category $\text{Glue}(R \to S, f_1, \ldots, f_t)$ has
direct sums and cokernels. The functor $\text{Can}$ commutes with direct sums
and cokernels. From these observations we find the displayed map is bijective
by writing $N$ as a cokernel of a map between free $R$-modules.
We omit the details.
\end{proof}

\begin{lemma}
\label{lemma-H0-inverse}
Assume $\varphi : R \to S$ is a flat ring map and
$I = (f_1, \ldots, f_t) \subset R$ is an ideal such that
$R/I \to S/IS$ is an isomorphism. Then the functor $H^0$
is a left quasi-inverse to the functor $\text{Can}$ of
Remark \ref{remark-glueing-data}.
\end{lemma}

\begin{proof}
This is a reformulation of
Lemma \ref{lemma-recover-module-from-glueing-data}.
\end{proof}

\begin{lemma}
\label{lemma-exact}
Assume $\varphi : R \to S$ is a flat ring map and let
$I = (f_1, \ldots, f_t) \subset R$ be an ideal.
Then $\text{Glue}(R \to S, f_1, \ldots, f_t)$ is an abelian category, and
the functor $\text{Can}$ is exact and commutes with arbitrary colimits.
\end{lemma}

\begin{proof}
Given a morphism
$(\varphi', \varphi_i) :
(M', M_i, \alpha_i, \alpha_{ij})
\to
(N', N_i, \beta_i, \beta_{ij})$
of the category $\text{Glue}(R \to S, f_1, \ldots, f_t)$
we see that its kernel exists and is equal to the object
$(\Ker(\varphi'), \Ker(\varphi_i), \alpha_i, \alpha_{ij})$
and its cokernel exists and is equal to the object
$(\Coker(\varphi'), \Coker(\varphi_i), \beta_i, \beta_{ij})$.
This works because $R \to S$ is flat, hence taking kernels/cokernels
commutes with $- \otimes_R S$. Details omitted.
The exactness follows from the $R$-flatness of $R_{f_i}$ and $S$, while
commuting with colimits follows as tensor products commute with colimits.
\end{proof}

\begin{lemma}
\label{lemma-equivalence-I-unit}
Let $\varphi : R \to S$ be a flat ring map and $(f_1, \ldots, f_t) = R$.
Then $\text{Can}$ and $H^0$ are quasi-inverse equivalences of categories
$$
\text{Mod}_R = \text{Glue}(R \to S, f_1, \ldots, f_t)
$$
\end{lemma}

\begin{proof}
Consider an object $\mathbf{M} = (M', M_i, \alpha_i, \alpha_{ij})$
of $\text{Glue}(R \to S, f_1, \ldots, f_t)$. By
Algebra, Lemma \ref{algebra-lemma-glue-modules}
there exists a unique module $M$ and isomorphisms
$M_{f_i} \to M_i$ which recover the glueing data $\alpha_{ij}$.
Then both $M'$ and $M \otimes_R S$ are $S$-modules
which recover the modules $M_i \otimes_R S$ upon localizing at $f_i$.
Whence there is a canonical isomorphism $M \otimes_R S \to M'$.
This shows that $\mathbf{M}$ is in the essential image of $\text{Can}$.
Combined with
Lemma \ref{lemma-H0-inverse}
the lemma follows.
\end{proof}

\begin{lemma}
\label{lemma-base-change-glue}
Let $\varphi : R \to S$ be a flat ring map and $I = (f_1, \ldots, f_t)$
and ideal. Let $R \to R'$ be a flat ring map, and set $S' = S \otimes_R R'$.
Then we obtain a commutative diagram of categories and functors
$$
\xymatrix{
\text{Mod}_R \ar[r]_-{\text{Can}} \ar[d]_{-\otimes_R R'} &
\text{Glue}(R \to S, f_1, \ldots, f_t) \ar[r]_-{H^0} \ar[d]^{-\otimes_R R'} &
\text{Mod}_R \ar[d]^{-\otimes_R R'} \\
\text{Mod}_{R'} \ar[r]^-{\text{Can}} &
\text{Glue}(R' \to S', f_1, \ldots, f_t) \ar[r]^-{H^0} &
\text{Mod}_{R'}
}
$$
\end{lemma}

\begin{proof}
Omitted.
\end{proof}

\begin{proposition}
\label{proposition-equivalence}
Assume $\varphi : R \to S$ is a flat ring map and
$I = (f_1, \ldots, f_t) \subset R$ is an ideal such that
$R/I \to S/IS$ is an isomorphism. Then $\text{Can}$ and
$H^0$ are quasi-inverse equivalences of categories
$$
\text{Mod}_R = \text{Glue}(R \to S, f_1, \ldots, f_t)
$$
\end{proposition}

\begin{proof}
We have already seen that $H^0 \circ \text{Can}$ is isomorphic to the
identity functor, see
Lemma \ref{lemma-H0-inverse}.
Consider an object $\mathbf{M} = (M', M_i, \alpha_i, \alpha_{ij})$
of $\text{Glue}(R \to S, f_1, \ldots, f_t)$.
We get a natural morphism
$$
\Psi :
(H^0(\mathbf{M}) \otimes_R S, H^0(\mathbf{M})_{f_i},
\text{can}_i, \text{can}_{ij})
\longrightarrow
(M', M_i, \alpha_i, \alpha_{ij}).
$$
Namely, by definition $H^0(\mathbf{M})$ comes equipped with compatible
$R$-module maps $H^0(\mathbf{M}) \to M'$ and $H^0(\mathbf{M}) \to M_i$.
We have to show that this map is an isomorphism.

\medskip\noindent
Pick an index $i$ and set $R' = R_{f_i}$. Combining
Lemmas \ref{lemma-base-change-glue} and \ref{lemma-equivalence-I-unit}
we see that $\Psi \otimes_R R'$ is an isomorphism.
Hence the kernel, resp.\ cokernel of $\Psi$ is a system of the form
$(K, 0, 0, 0)$, resp.\ $(Q, 0, 0, 0)$. Note that
$H^0((K, 0, 0, 0)) = K$, that $H^0$ is left exact, and that by
construction $H^0(\Psi)$ is bijective. Hence we see $K = 0$, i.e.,
the kernel of $\Psi$ is zero.

\medskip\noindent
The conclusion of the above is that we obtain a short exact sequence
$$
0 \to H^0(\mathbf{M}) \otimes_R S \to M' \to Q \to 0
$$
and that $M_i = H^0(\mathbf{M})_{f_i}$.
Note that we may think of $Q$ as an $R$-module which is $I$-power
torsion so that $Q = Q \otimes_R S$. By
Lemma \ref{lemma-neighbourhood-extensions}
we see that there exists a commutative diagram
$$
\xymatrix{
0 \ar[r] &
H^0(\mathbf{M}) \ar[r] \ar[d] &
E \ar[r] \ar[d] &
Q \ar[r] \ar[d] &
0 \\
0 \ar[r] &
H^0(\mathbf{M}) \otimes_R S \ar[r] &
M' \ar[r] &
Q \ar[r] &
0
}
$$
with exact rows. This clearly determines an isomorphism
$\text{Can}(E) \to (M', M_i, \alpha_i, \alpha_{ij})$
in the category $\text{Glue}(R \to S, f_1, \ldots, f_t)$
and we win. (Of course, a posteriori we have $Q = 0$.)
\end{proof}

\begin{lemma}
\label{lemma-application-formal-glueing}
Let $\varphi : R \to S$ be a flat ring map and let $I \subset R$ be a
finitely generated ideal such that $R/I \to S/IS$ is an isomorphism. 
\begin{enumerate}
\item Given an $R$-module $N$, an $S$-module $M'$ and an $S$-module
map $\varphi : M' \to N \otimes_R S$ whose kernel and cokernel are
$I$-power torsion, there exists an $R$-module map
$\psi : M \to N$ and an isomorphism $M \otimes_R S = M'$
compatible with $\varphi$ and $\psi$.
\item Given an $R$-module $M$, an $S$-module $N'$ and an $S$-module
map $\varphi : M \otimes_R S \to N'$ whose kernel and cokernel are
$I$-power torsion, there exists an $R$-module map
$\psi : M \to N$ and an isomorphism $N \otimes_R S = N'$
compatible with $\varphi$ and $\psi$.
\end{enumerate}
In both cases we have $\Ker(\varphi) \cong \Ker(\psi)$ and
$\Coker(\varphi) \cong \Coker(\psi)$.
\end{lemma}

\begin{proof}
Proof of (1). Say $I = (f_1, \ldots, f_t)$. It is clear that
the localization $\varphi_{f_i}$ is an isomorphism.
Thus we see that $(M', N_{f_i}, \varphi_{f_i}, can_{ij})$ is an
object of $\text{Glue}(R \to S, f_1, \ldots, f_t)$, see
Remark \ref{remark-glueing-data}.
By Proposition \ref{proposition-equivalence}
we conclude that there exists an $R$-module $M$ such that
$M' = M \otimes_R S$ and $N_{f_i} = M_{f_i}$ compatibly
with the isomorphisms $\varphi_{f_i}$ and $can_{ij}$. There is a
morphism
$$
(M \otimes_R S, M_{f_i}, can_i, can_{ij}) =
(M', N_{f_i}, \varphi_{f_i}, can_{ij})
\to
(N \otimes_R S, N_{f_i}, can_i, can_{ij})
$$
of $\text{Glue}(R \to S, f_1, \ldots, f_t)$
which uses $\varphi$ in the first component. This
corresponds to an $R$-module map $\psi : M \to N$ (by the equivalence of
categories of Proposition \ref{proposition-equivalence}).
The composition of the base change of $M \to N$ with the
isomorphism $M' \cong M \otimes_R S$ is $\varphi$, in other words
$M \to N$ is compatible with $\varphi$.

\medskip\noindent
Proof of (2). This is just the dual of the argument above.
Namely, the localization $\varphi_{f_i}$ is an isomorphism.
Thus we see that $(N', M_{f_i}, \varphi_{f_i}^{-1}, can_{ij})$ is an
object of $\text{Glue}(R \to S, f_1, \ldots, f_t)$, see
Remark \ref{remark-glueing-data}.
By Proposition \ref{proposition-equivalence}
we conclude that there exists an $R$-module $N$ such that
$N' = N \otimes_R S$ and $N_{f_i} = M_{f_i}$ compatibly
with the isomorphisms $\varphi_{f_i}^{-1}$ and $can_{ij}$. There is a
morphism
$$
(M \otimes_R S, M_{f_i}, can_i, can_{ij}) \to
(N', M_{f_i}, \varphi_{f_i}, can_{ij}) =
(N \otimes_R S, N_{f_i}, can_i, can_{ij})
$$
of $\text{Glue}(R \to S, f_1, \ldots, f_t)$
which uses $\varphi$ in the first component. This
corresponds to an $R$-module map $\psi : M \to N$ (by the equivalence of
categories of Proposition \ref{proposition-equivalence}).
The composition of the base change of $M \to N$ with the
isomorphism $N' \cong N \otimes_R S$ is $\varphi$, in other words
$M \to N$ is compatible with $\varphi$.

\medskip\noindent
The final statement follows for example from
Lemma \ref{lemma-neighbourhood-equivalence}.
\end{proof}

\noindent
Next, we specialize Proposition \ref{proposition-equivalence} to get something
more usable. Namely, if $I = (f)$ is a principal ideal then the objects
of $\text{Glue}(R \to S, f)$ are simply triples $(M', M_1, \alpha_1)$
and there is {\it no} cocycle condition to check!

\begin{theorem}
\label{theorem-formal-glueing}
Let $R$ be a ring, and let $f \in R$.
Let $\varphi : R \to S$ be a flat ring map inducing an isomorphism
$R/fR \to S/fS$. Then the functor
$$
\text{Mod}_R
\longrightarrow
\text{Mod}_S \times_{\text{Mod}_{S_f}} \text{Mod}_{R_f},
\quad
M
\longmapsto
(M \otimes_R S, M_f, \text{can})
$$
is an equivalence.
\end{theorem}

\begin{proof}
The category appearing on the right side of the arrow
is the category of triples $(M', M_1, \alpha_1)$ where $M'$ is an
$S$-module, $M_1$ is a $R_f$-module, and
$\alpha_1 : M'_f \to M_1 \otimes_R S$ is a $S_f$-isomorphism, see
Categories, Example \ref{categories-example-2-fibre-product-categories}.
Hence this theorem is a special case of
Proposition \ref{proposition-equivalence}.
\end{proof}

\noindent
A useful special case of
Theorem \ref{theorem-formal-glueing}
is when $R$ is Noetherian, and $S$ is a completion of $R$ at an
element $f$. The completion $R \to S$ is flat, and the functor
$M \mapsto M \otimes_R S$ can be identified with the $f$-adic
completion functor when $M$ is finitely generated. To state
this more precisely, let $\text{Mod}^{fg}_R$ denote the category
of finitely generated $R$-modules.

\begin{proposition}
\label{proposition-formal-glueing}
Let $R$ be a Noetherian ring.
Let $f \in R$ be an element.
Let $R^\wedge$ be the $f$-adic completion of $R$.
Then the functor $M \mapsto (M^\wedge, M_f, \text{can})$
defines an equivalence
$$
\text{Mod}^{fg}_R
\longrightarrow
\text{Mod}^{fg}_{R^\wedge}
\times_{\text{Mod}^{fg}_{(R^\wedge)_f}}
\text{Mod}^{fg}_{R_f}
$$
\end{proposition}

\begin{proof}
The ring map $R \to R^\wedge$ is flat by
Algebra, Lemma \ref{algebra-lemma-completion-flat}.
It is clear that $R/fR = R^\wedge/fR^\wedge$.
By
Algebra, Lemma \ref{algebra-lemma-completion-tensor}
the completion of a finite $R$-module $M$ is equal to $M \otimes_R R^\wedge$.
Hence the displayed functor of the proposition is equal to the
functor occurring in
Theorem \ref{theorem-formal-glueing}.
In particular it is fully faithful. Let $(M_1, M_2, \psi)$ be an
object of the right hand side. By
Theorem \ref{theorem-formal-glueing}
there exists an $R$-module $M$ such that
$M_1 = M \otimes_R R^\wedge$ and $M_2 = M_f$. As $R \to R^\wedge \times R_f$
is faithfully flat we conclude from
Algebra, Lemma \ref{algebra-lemma-cover}
that $M$ is finitely generated, i.e., $M \in \text{Mod}^{fg}_R$.
This proves the proposition.
\end{proof}

\begin{remark}
\label{remark-formal-glueing-algebras}
The equivalences of
Proposition \ref{proposition-equivalence},
Theorem \ref{theorem-formal-glueing}, and
Proposition \ref{proposition-formal-glueing}
preserve properties of modules. For example if
$M$ corresponds to $\mathbf{M} = (M', M_i, \alpha_i, \alpha_{ij})$
then $M$ is finite, or finitely presented, or flat, or projective over $R$
if and only if $M'$ and $M_i$ have the corresponding property
over $S$ and $R_{f_i}$. This follows from the fact that
$R \to S \times \prod R_{f_i}$ is faithfully flat and
descend and ascent of these properties along faithfully flat maps, see
Algebra, Lemma \ref{algebra-lemma-descend-properties-modules} and
Theorem \ref{algebra-theorem-ffdescent-projectivity}.
These functors also preserve the $\otimes$-structures on either side.
Thus, it defines equivalences of various categories
built out of the pair $(\text{Mod}_R, \otimes)$, such as the category of
algebras.
\end{remark}

\begin{remark}
\label{remark-topological-analogue}
Given a differential manifold $X$ with a compact closed submanifold $Z$
having complement $U$, specifying a sheaf on $X$ is the same as specifying
a sheaf on $U$, a sheaf on an unspecified tubular neighbourhood $T$ of $Z$ in
$X$, and an isomorphism between the two resulting sheaves along $T \cap U$.
Tubular neighbourhoods do not exist in algebraic geometry as such, but
results such as
Proposition \ref{proposition-equivalence},
Theorem \ref{theorem-formal-glueing}, and
Proposition \ref{proposition-formal-glueing}
allow us to work with formal neighbourhoods instead.
\end{remark}





